\documentclass{article}
\usepackage{graphicx}

%PREAMBLE
\title{Activities 2016 / 2017}
\author{Nora Juhasz}

%CONTENT
\begin{document}

\maketitle
\newpage

\begin{abstract}

This document describes the university classes, GSSI classes, summer schools and seminars attended in the 2016 / 2017 academic year. The bibliography contains the list of the books and articles that mean the basis of this research project.\\

\end{abstract}
\newpage

\tableofcontents
\newpage

%UNIVAQ
\section{University of L'Aquila}

\subsection{Numerical Analysis 2 - Prof. Guglielmi}
Euler method. Steepest descent gradient method. Conjugate directions method. Matlab.

\subsection{High Performance Parallel Computing - Prof. Shcherbatyy, Prof. Pera}
Basic Linux commands. 

%GSSI
\section{Gran Sasso Science Institute}

\subsection{Introduction to Finite Elements Method - Prof. Boffi}
Finite differences, finite elements.

\subsection{Decay property for hyperbolic systems with dissipation - Prof. Kawashima}
The aim of the course is to study dissipative structure and decay property for hyperbolic systems with dissipation. Consider the following two equations.

\begin{equation} \label{eq:symmetricdiffusion}
u_{t} + Au_{x} = Bu_{xx}
\end{equation}

\begin{equation} \label{eq:symmetricrelaxation}
u_{t} + Au_{x} -Lu = 0
\end{equation}

Equation \ref{eq:symmetricdiffusion} describes the symmetric diffusion process while symmetric relaxation is modeled by equation \ref{eq:symmetricrelaxation}.

It is always assumed that the matrix $A$ is symmetric. Throughout the course of this lecture three different versions of these two equations are examined in terms of the assumptions that are made on matrices $B$ and $L$, and what type of additional effects are involved. The simplest case is when one assumes $B$ and $L$ to be symmetric and nonnegative. From this base case we can move in two different directions and arrive at two different versions of these systems. In the first variation, the very same equations are considered with modified conditions, ie. matrices $B$ and $L$ are still assumed to be non-negative, but the symmetry assumption is dropped. In the second variation one considers the original assumptions on symmetry and non-negativity, but the equations are modified with a memory part. More specifically, the equations in the third case are the following.

\begin{equation} \label{eq:symmetricdiffusionmemo}
u_{t} + Au_{x} = a(t) \star Bu_{xx}
\end{equation}

\begin{equation} \label{eq:symmetricrelaxationmemo}
u_{t} + Au_{x} -a(t) \star Lu = 0
\end{equation}

Fourier transform.

\subsection{Advanced Topics on Numerical Methods - Prof. Guglielmi}
A-stability, B-stability, Runge-Kutta.

\subsection{Introduction to the Incompressible Navier-Stokes equation and the mathematical theory of Turbulence - Prof. Marcati}
Euler and Navier-Stokes equations.

%SEMINARS
\section{Seminars}

\subsection{PDEs Day, 21st March 2017}

\subsubsection{Viscous vortex rings with axial symmetry - Prof. Gallay}

\subsubsection{A two phase model with cross and self attractive interactions - Prof. Novaga}

\subsubsection{Traveling waves for collective movements on a network - Prof. Corli}

%SUMMER SCHOOLS
\section{Summer schools}

\nocite{*}
\bibliographystyle{apalike}
\bibliography{bibliography}

\end{document}
